% ECONOMICS_PROBLEM: {"id": "G51-644", "title": "Long-run effects of defaults", "jel_code": "G51", "source_id": "OP-0065"}
% !TeX program = xelatex
\documentclass[11pt,a4paper]{article}
\usepackage[margin=23mm]{geometry}
\usepackage{fontspec}
\setmainfont{Latin Modern Roman}
\usepackage{amsmath,amssymb,mathtools}
\usepackage{xcolor,enumitem,booktabs,longtable,array,xurl,hyperref}
\definecolor{muted}{HTML}{53636C}
\hypersetup{colorlinks=true,urlcolor=blue,linkcolor=blue}
\setlength{\parindent}{0pt}
\setlength{\parskip}{5pt}
\newcommand{\R}{\mathbb R}
\newcommand{\E}{\mathbb E}
\newcommand{\N}{\mathbb N}
\newcommand{\ind}{\mathbf 1}
\newcommand{\Var}{\operatorname{Var}}
\newcommand{\Cov}{\operatorname{Cov}}
\newcommand{\supp}{\operatorname{supp}}
\DeclareMathOperator*{\argmin}{arg\,min}
\DeclareMathOperator*{\argmax}{arg\,max}
\newcommand{\entrymeta}[3]{{\sffamily\small\color{muted}\textbf{\texttt{#1}}\quad|\quad #2\par #3\par}}
\newcommand{\Problemref}[1]{\texttt{#1}}
\newcommand{\JELref}[2]{\texttt{#2} (JEL \texttt{#1})}
\begin{document}
\section*{Long-run effects of defaults}
\textbf{Problem ID:} \texttt{G51-644}\quad \textbf{JEL:} \texttt{G51}\par
% BEGIN ECONOMICS STATEMENT
\label{problem:OP-0065}
\label{atlas:prob:B5}

\noindent\textbf{Original identifier:} B5 (atlas item 65). \textbf{Classification:} Longitudinal causal and welfare research program. \textbf{Original tractability label:} Medium (editorial, not a complexity result).

\subsection*{Definitions and assumptions}

Let $D\in\{0,1\}$ be a randomly assigned or conditionally exogenous enrollment default. For horizon $t=0,\ldots,T$, define contribution $C_t(d)$, retirement-plan wealth $R_t(d)$, other assets $O_t(d)$, debt $L_t(d)$, and net wealth $N_t(d)=R_t(d)+O_t(d)-L_t(d)$. Specify taxes, employer contributions, withdrawals, employment exits, and consumption measurement. An intention-to-treat effect averages over the baseline eligible population, including later leavers. For welfare, supply an evaluator $w$ and a finite-horizon utility functional $U^w$; do not infer this functional solely from observed adherence to the default.

\subsection*{Formal research question}

Identify the trajectories $\Delta_C(t)=\mathbb E[C_t(1)-C_t(0)]$ and $\Delta_N(t)=\mathbb E[N_t(1)-N_t(0)]$, and the welfare difference
\[
\Delta_W^w(T)=\mathbb E[U^w(c_{0:T}(1),\ell_{0:T}(1))-U^w(c_{0:T}(0),\ell_{0:T}(0))],
\]
where $c$ and $\ell$ denote consumption and leisure. Determine whether economically meaningful persistence, $|\Delta_N(t)|\geq\varepsilon$ for all observed $t\in\{t_0,\ldots,T\}$, holds for prespecified $\varepsilon>0$. Derive bounds under missing post-separation outcomes and model sensitivity for extrapolation beyond $T$.

\subsection*{Interpretation and scope}

Finite follow-up can establish persistence over a specified horizon, not a permanent effect over an infinite lifetime. The relevant savings outcome includes borrowing and offsetting asset changes; a rise in plan balances alone does not identify net saving. Post-treatment employment status and plan participation must not be conditioned upon without accounting for selection. A randomized default identifies total policy effects but generally does not separately identify inertia, learning, and preference change. Mechanism claims need additional interventions, such as removing the default, changing information, or varying switching costs, with explicit exclusion restrictions. Welfare can improve or deteriorate for people induced to save, depending on liquidity needs and the selected evaluator. Long-run data linkage should track withdrawals, replacement jobs, and debt, and handle mortality as part of a defined outcome rather than silently treating survivors as the original population.

\subsection*{Source audit and status}

[S1]--[S3] verify the original default literature; Beshears et al. (2009) is resolved contextually to the retirement-default chapter, since author-year alone is not unique. [S4], a 2024 working paper, specifically qualifies naive extrapolation of initial effects through medium- and long-run dynamics. Thus persistence is an active empirical question with existing answers in particular settings, not a universal unanswered proposition. The net-wealth and criterion-specific welfare extension above is editorial.

\subsection*{References}

\begin{enumerate}[label={[S\arabic*]},leftmargin=*]

\item Brigitte C. Madrian and Dennis F. Shea (2001). \emph{The Power of Suggestion: Inertia in 401(k) Participation and Savings Behavior}. Quarterly Journal of Economics 116(4), 1149--1187. \url{https://www.nber.org/papers/w7682}\par \textit{Source role:} Evidence on automatic enrollment; NBER record links published version. \textit{Locator:} Abstract and article metadata.

\item Gabriel D. Carroll, James J. Choi, David Laibson, Brigitte C. Madrian, and Andrew Metrick (2009). \emph{Optimal Defaults and Active Decisions}. Quarterly Journal of Economics 124(4), 1639--1674. \url{https://doi.org/10.1162/qjec.2009.124.4.1639}\par \textit{Source role:} Model and evidence for alternative enrollment regimes. \textit{Locator:} Abstract and article metadata.

\item John Beshears, James J. Choi, David Laibson, and Brigitte C. Madrian (2009). \emph{The Importance of Default Options for Retirement Saving Outcomes: Evidence from the United States}. In Social Security Policy in a Changing Environment, University of Chicago Press, 167--195. \url{https://www.nber.org/books-and-chapters/social-security-policy-changing-environment/importance-default-options-retirement-saving-outcomes-evidence-united-states}\par \textit{Source role:} Review of retirement default evidence. \textit{Locator:} Abstract and article metadata.

\item James J. Choi, David Laibson, Jordan Cammarota, Richard Lombardo, and John Beshears (2024). \emph{Smaller than We Thought? The Effect of Automatic Savings Policies}. NBER Working Paper 32828. \url{https://www.nber.org/papers/w32828}\par \textit{Source role:} Recent qualification of long-run default effects. \textit{Locator:} Abstract; medium- and long-run dynamics.

\end{enumerate}

\subsection*{Common conventions: Source mathematical conventions}

Unless an entry states otherwise, agent and firm sets are finite. State and action spaces are measurable, strategies and outcome maps are measurable, probability laws are specified, and expectations exist for the targets under discussion. Infinite-horizon discount factors lie in $(0,1)$ and discounted objectives are finite. Norms, tolerances and complexity input sizes must be specified for computational comparisons. Constants or primitives that appear locally are inputs, not empirical facts asserted by this catalogue.

\textbf{Identification.} A statistical model $m\in\mathcal M$ implies an observable law $P_m$ and target $\theta(m)$. At law $P$, its identified set is
\[
\Theta(P;\mathcal M)=\{\theta(m):m\in\mathcal M,\ P_m=P\}.
\]
Point identification holds when this set is a singleton, or equivalently when $P_{m_1}=P_{m_2}$ implies $\theta(m_1)=\theta(m_2)$ for all compatible models. Sharp bounds mean bounds attained, or approached when the boundary is not attained, by compatible models. Writing this set is a definition, not a solution: the substantive task is to establish informative restrictions and derive or estimate the set. An unrestricted-confounding model may give only trivial bounds.

\textbf{Inference.} Identification, estimability and valid uncertainty quantification are separate questions. Fix $\alpha\in(0,1)$. A confidence procedure $C_n$ over a declared law class $\mathcal P$ has asymptotic uniform coverage at level $1-\alpha$ when
\[
\liminf_{n\to\infty}\inf_{P\in\mathcal P}\Pr_P\{\theta(P)\in C_n\}\geq1-\alpha.
\]
For set-identified targets the coverage event must be defined separately, for example coverage of the full identified set. If an entry uses identification-indexed models instead of uniquely indexed laws, the same coverage convention is applied over those models. Pointwise limits do not establish uniform validity, and asymptotic validity does not guarantee finite-sample precision. Sample sizes, dependence structures and weak-identification sequences are part of each application.

\textbf{Counterfactuals.} $Y_i(a)$ is a potential outcome under a specified intervention, including its timing, implementation and versions. It is not $Y_i$ conditional on voluntarily choosing $a$. A full intervention vector permits interference; shorthand scalar treatment notation does not silently impose no interference. Assignment, selection, attrition, anticipation and measurement must be specified. A claim about an instrument's local effect is not automatically a population policy effect.

\textbf{Economic equilibrium.} A counterfactual model includes best responses, constraints, feasibility, market clearing, state transitions and expectations consistent with the stated information. When several equilibria exist, either a declared selector is an input or the target is set-valued. Comparisons across policy regimes must use the stated selection convention. Reduced-form relationships are not presumed invariant after an equilibrium-changing intervention.

\textbf{Optimization and welfare.} Optimization is over the stated feasible set. If attainment is not established, $\sup$ or $\inf$ and $\varepsilon$-optimal choices replace an asserted maximizer or minimizer. For example, with value $V=\sup_{a\in\mathcal A}W(a)<\infty$, an $\varepsilon$-optimal set is $\{a\in\mathcal A:W(a)\geq V-\varepsilon\}$ for $\varepsilon>0$. Benefits, costs, risk and health or environmental outcomes can be aggregated only using declared units and conversion weights. Interpersonal utility weights, the treatment of fiscal transfers and foreign profits, discounting, and the behavioral welfare criterion are explicit inputs. A comparative-static derivative additionally requires existence and appropriate differentiability of the selected counterfactual.

\textbf{Sources.} Local labels [S1], [S2], and so forth restart in every question. Locators identify the relevant abstract, model, section or result at the level actually checked; a bibliographic or abstract-level verification is not represented as inspection of an entire proof. Working papers are distinguished from later publications and changing versions. Source summaries are paraphrased. All URL checks and editorial formulations refer to the audit date stated above.
% END ECONOMICS STATEMENT
\end{document}
